% ECONOMICS_PROBLEM: {"id": "D63-72", "title": "EFM with nonadditive mixed items and desirable cake", "jel_code": "D63", "source_id": "OP-0594"}
% !TeX program = xelatex
\documentclass[11pt,a4paper]{article}
\usepackage[margin=23mm]{geometry}
\usepackage{fontspec}
\setmainfont{Latin Modern Roman}
\usepackage{amsmath,amssymb,mathtools}
\usepackage{xcolor,enumitem,booktabs,longtable,array,xurl,hyperref}
\definecolor{muted}{HTML}{53636C}
\hypersetup{colorlinks=true,urlcolor=blue,linkcolor=blue}
\setlength{\parindent}{0pt}
\setlength{\parskip}{5pt}
\newcommand{\R}{\mathbb R}
\newcommand{\E}{\mathbb E}
\newcommand{\N}{\mathbb N}
\newcommand{\ind}{\mathbf 1}
\newcommand{\Var}{\operatorname{Var}}
\newcommand{\Cov}{\operatorname{Cov}}
\newcommand{\supp}{\operatorname{supp}}
\DeclareMathOperator*{\argmin}{arg\,min}
\DeclareMathOperator*{\argmax}{arg\,max}
\newcommand{\entrymeta}[3]{{\sffamily\small\color{muted}\textbf{\texttt{#1}}\quad|\quad #2\par #3\par}}
\newcommand{\Problemref}[1]{\texttt{#1}}
\newcommand{\JELref}[2]{\texttt{#2} (JEL \texttt{#1})}
\begin{document}
\section*{EFM with nonadditive mixed items and desirable cake}
\textbf{Problem ID:} \texttt{D63-72}\quad \textbf{JEL:} \texttt{D63}\par
% BEGIN ECONOMICS STATEMENT
\label{problem:OP-0594}
% ECONOMICS_PROBLEM: {"local_id": "ECON-008", "title": "EFM with nonadditive mixed items and desirable cake", "kind": "P", "audit_date": "2026-10-08", "source_identity_checked": true, "agenda_or_question_support_checked": true, "source_role": "see per-entry audit and locators", "current_open_certified": false}


\label{audited:ECON-008}

{\sffamily\small\color{muted}\textbf{ECON-008} | P | Explicit published open question, at source date\par}

\subsection*{Definitions and formal target}

Let $N=\{1,\ldots,n\}$, $M$ be finite, and $v_i:2^M\to\mathbb R$ satisfy $v_i(\varnothing)=0$. Doubly monotone means there is a partition $M=G_i\sqcup B_i$ such that for every $S\subseteq M$ and $g\notin S$,
\[
g\in G_i\Rightarrow v_i(S\cup\{g\})\geq v_i(S),\qquad
g\in B_i\Rightarrow v_i(S\cup\{g\})\leq v_i(S).
\]
The sign partition may depend on $i$; an item with identically zero marginal contribution may be assigned to either part. Cake is $C=[0,1]$ with nonnegative Lebesgue-integrable densities $f_i$ and $V_i(D)=\int_Df_i(x)\,dx$. A complete allocation partitions the items into $O_i$ and cake into measurable pieces $D_i$, up to null endpoints. Utilities are $u_i(O,D)=v_i(O)+V_i(D)$.

Use the \emph{stronger} EFM definition in [S1], Definition 2.3. For every ordered pair $i,j$, require either $u_i(O_i,D_i)\geq u_i(O_j,D_j)$, or
\[
 V_i(D_j)=0\quad\text{and}\quad
 \exists g\in O_i\cup O_j:\quad
 v_i(O_i\setminus\{g\})\geq v_i(O_j\setminus\{g\}).
\]
Does every instance in this domain have a complete EFM allocation?

\subsection*{Clarifications and required assumptions}

The fallback compares \emph{item valuations only}. The draft instead included $V_i(D_i)$ in that comparison and therefore stated a weaker condition. This has been corrected. An item in $O_i$ is removed from the own bundle and an item in $O_j$ from the compared bundle; removal of the same named item from the other disjoint bundle does nothing. With empty item bundles, cake envy must be ruled out by the first branch. These density-induced measures are absolutely continuous and atomless; a general atomless measure need not have a Lebesgue density.

\subsection*{Variants and changes of scope}

\textbf{Weak EFM.} Keep $V_i(D_i)$ in the fallback inequality; because $V_i(D_j)=0$, the fallback becomes $v_i(O_i\setminus\{g\})+V_i(D_i)\geq v_i(O_j\setminus\{g\})$. This is weaker than the main definition. \textbf{Indivisible-only compatibility.} Without cake, ask for EF1 and an envy-freeable allocation: some nonnegative vector $p$ must satisfy $v_i(O_i)+p_i\geq v_i(O_j)+p_j$ for every pair. \textbf{Additive items.} Additivity is a settled subcase in [S1]; it must not be relisted as the main open question. \textbf{Divisible chore.} Negative cake disutility changes both the model and fairness direction. \textbf{Measure generalization.} Replace density-induced cake values by arbitrary finite nonnegative atomless Borel measures, meaning $V_i(\{x\})=0$ for every point $x$. This enlarges the domain and is an editorial extension, not silently the source's stated model.

\subsection*{Audit conclusion and source relationship}

[S1], Section 5, explicitly poses the doubly-monotone extension and the indivisible-only EF1/envy-freeability question. Definition 2.3 and Remark 2.4 reveal the draft's mismatch. The revised main condition follows that paper; weak EFM is now a named variant.

\subsection*{References and exact source roles}

{\footnotesize\raggedright
\par\noindent\textbf{[S1]} Haris Aziz, Xinhang Lu, Simon Mackenzie, and Mashbat Suzuki (2025). \emph{Fair Division with Indivisible Goods, Chores, and Cake}. arXiv:2511.04891v1, 7 November 2025; preprint.\par \textit{Verified locator / source role:} Section 5, Discussion: doubly-monotone extension; definitions of EFM.\par \url{https://arxiv.org/html/2511.04891v1}\par\smallskip
}

\subsection*{Common conventions: Additional-source status and mathematical conventions}

\label{audited:conventions}

Of the 100 supplied ECON entries, 65 distinct problems appear here; 32 close
matches refine earlier entries, and three duplicate questions map to existing
entries. Appendix~\ref{app:audited-crosswalk} lists every destination.
The attachment's P/R/S/U distinctions and source audits are retained. Its
precise mathematical formalizations are editorial unless the individual audit
says otherwise. Candidate subsidy targets ECON-004--006 retain their U flags;
the resolved approval-core question ECON-007 updates OP-0202 above.
The following supplied conventions govern both this part and the attributed
ECON refinements elsewhere in the catalogue, subject to local restrictions.

\textbf{Parameterized questions.} A problem family lists inputs that must be fixed before an instance is evaluated: population, horizon, policies, information, data law, model class and welfare weights. These are required choices, not quantities the citations are assumed to specify. Undefined applications should not be treated as identified estimands.

\textbf{Indices and integrability.} Sets of agents, firms and regions are finite unless a population measure is explicitly used. \(i,j\) index units, \(r\) regions, \(t\) dates and \(h\) horizons. \(\mathbb E\) and \(\Pr\) refer to the declared population and law; \(\mathbf1\{A\}\) is an indicator. Every displayed expectation and discounted sum needs existence in the stated domain. Infinite horizons use a discount in \((0,1)\) and a convergence condition; finite horizons may use discount one.

\textbf{Optimization and existence.} A displayed \(\max\), \(\min\), or \(\arg\max\) is conditional on attainment. For finite feasible menus this is immediate; general spaces need continuity and compactness/coercivity or another existence argument. Without attainment, the target is the supremum/infimum and arbitrarily near-optimal feasible choices. \(\inf\varnothing=+\infty\). Empty feasibility and an unidentified target are different issues.

\textbf{Potential outcomes.} \(Y_i(z)\) is the outcome under a specified intervention, not the observed conditional mean among people choosing \(z\). In binary comparisons \(\Delta Y_i=Y_i(1)-Y_i(0)\). Specify assignment, treatment versions, compliance, information and observation horizon. Interference is allowed under a full policy vector; compressed exposure notation needs a justified mapping. No interference, consistency and overlap are substantive assumptions, not automatic consequences of notation.

\textbf{Populations and observation.} Fix baseline eligibility and population weights. Conditioning on post-policy employment, survival, relocation or program uptake can change the population being compared. Death is not ordinary loss to follow-up; outcomes truncated by death need a composite convention, joint survival target, or explicitly defined survivor stratum. Fertility-changing interventions need a population functional rather than an unexplained fixed descendant cohort. Measurement and missingness models must be supplied.

\textbf{Identification.} A structural model \(m\in\mathcal M\) implies observable law \(P_m\) and target \(\theta(m)\). Identification means
\[
 P_{m_1}=P_{m_2}\ \Longrightarrow\ \theta(m_1)=\theta(m_2)
 \quad\text{for all }m_1,m_2\in\mathcal M .
\]
Without identification, the target is
\[
 \Theta(P;\mathcal M)=\{\theta(m):m\in\mathcal M,\ P_m=P\}.
\]
Writing \(\theta(P)\) before establishing identification can hide incompatible latent models. Confidence sets additionally need a sampling law, confidence level, asymptotic sequence and uniformity class. Point identification, coverage and informative precision are separate properties.

\textbf{Mediation.} Total effects of randomized assignment are distinct from cross-world natural indirect effects. Belief/effort-channel effects require a meaningful mediator intervention and additional measurement, exclusion or mediation assumptions. Randomization of the initial policy alone does not supply those assumptions. Report bounds or interventional alternatives when the stronger target is unsupported.

\textbf{Equilibrium.} A counterfactual \(W(z)\), \(p(z)\), or \(\mu_t^z\) requires individual optimization, feasibility, government/resource budgets, market clearing and state-transition laws. State equilibrium existence and selection, or report ranges over compatible equilibria. Initial-state integration includes subsequent endogenous transitions; current-state integration must use the policy-dependent distribution. Reduced-form invariance after policy is not assumed.

\textbf{Welfare accounts.} Scores, utility, health, dollars and ecological quantities require explicit conversion before addition. Fees, taxes, subsidies, wages and extortion payments are transfers between parties; distributional weights can make them welfare relevant, but their face value is not automatically a resource loss. State ownership, geographic boundaries, foreign-income weights and fiscal finance. Do not add profits to household welfare already including dividends, or count land capitalization and the underlying benefit twice. A benefit-cost expression must distinguish gross benefits from net welfare and subtract every real cost exactly once.

\textbf{Derivatives and risk.} Log derivatives require positive arguments; local responses require admissible perturbations and specified equilibrium branches. Differentiation under expectations needs regularity/domination, and boundaries may allow only one-sided derivatives. Risk uses a specified law; ambiguity uses a declared nonempty law set on a common history space. Adaptive policies must be nonanticipative. A robust criterion is an editorial choice unless attributed explicitly.

\textbf{Scope overlaps.} Allocation, collective choice, contracts, household bargaining and coordinated policy overlap economics with optimization and game theory. They are retained because they appeared in the original catalogue; no claim of a strictly game-theory-free selection is made.
% END ECONOMICS STATEMENT
\end{document}
